\chapter{Protokół badania}
\label{appendix:protokol}

\section*{Informacje ogólne}

Niniejszy dokument opisuje procedurę sesji badawczej przeprowadzonej 
w ramach pracy inżynierskiej pt. \textit{„Sterowanie kursorem myszy 
za pomocą potencjałów wywoływanych stanu ustalonego"}.

\begin{itemize}
    \item \textbf{Czas trwania sesji:} około 20--30 minut
    \item \textbf{Wymagania sprzętowe:} komputer z monitorem LCD 
    (60~Hz), zestaw EEG OpenBCI Cyton, elektrody żelowe Ag/AgCl
    \item \textbf{Miejsce:} pomieszczenie o kontrolowanym oświetleniu
\end{itemize}

\section*{Przebieg sesji}

\paragraph{Krok 1 — Przygotowanie uczestnika (5 minut)}
Uczestnik zapoznaje się z celem badania oraz zasadą działania 
interfejsu SSVEP. Prowadzący informuje, że zadanie polega na 
skupianiu wzroku na migających kwadratach widocznych na ekranie 
i że nie wymaga żadnych ruchów ciała.

\paragraph{Krok 2 — Montaż elektrod (5 minut)}
Elektrody EEG umieszczane są w punktach O1, O2 i Oz zgodnie 
z systemem 10-20. Impedancja elektrod sprawdzana jest przed 
rozpoczęciem pomiaru. Uczestnik siada w odległości około 
60--70~cm od monitora.

\paragraph{Krok 3 — Sesja zapoznawcza (3 minuty)}
Uruchamiany jest interfejs stymulujący. Uczestnik przez 
około 3 minuty swobodnie eksploruje działanie systemu, 
bez rejestrowania wyników. Celem jest oswojenie się 
z opóźnieniem systemu i sposobem skupiania wzroku.

\paragraph{Krok 4 — Zadanie 1: Klikanie w cele}
Na ekranie wyświetlane są kolejno trzy czerwone kwadraty 
o wymiarach $200\times200$, $100\times100$ oraz $50\times50$ 
pikseli. Zadaniem uczestnika jest przemieszczenie kursora 
i kliknięcie w każdy z nich przy użyciu wyłącznie interfejsu 
BCI. Zadanie powtarzane jest trzykrotnie. Mierzony jest czas 
ukończenia każdego podejścia. Prowadzący rejestruje czas 
przy użyciu stopera.

\paragraph{Krok 5 — Zadanie 2: Skuteczność komend}
Prowadzący wydaje ustnie 15 poleceń w kolejności losowej 
(po 3 powtórzenia każdej z 5 komend: góra, dół, lewo, prawo, 
kliknięcie). Po każdym poleceniu uczestnik ma 8 sekund na 
jego wykonanie. Prowadzący odnotowuje wynik (sukces / 
niepowodzenie / błędna komenda) na formularzu obserwacyjnym.

\paragraph{Krok 6 — Ankieta (5 minut)}
Uczestnik wypełnia ankietę oceny systemu 
(Załącznik~\ref{appendix:ankieta}).